\FloatBarrier
\section{本章小结}
\label{sec:differential-perturbation-summary}

本章用三个量回答输入扰动的影响。微分给出指定方向上的一阶响应$\mathrm dF_{\bm x}(\bm h)$；局部条件数是这些响应在全部允许方向上的最坏放大倍数，可微时等于微分在允许方向子空间上的算子范数；余项说明给定步长下线性预测还差多少。三者回答不同的问题：条件数只描述无穷小扰动，可以有限而映射不可微，也可以为零而有限变化不为零；余项界需要导数沿整条扰动线段的Lipschitz型控制，二阶导数的统一界是取得这一控制的充分办法。

显式矩阵运算与隐式谱分支遵循同一套步骤：确定输出类型与误差尺度，限定允许的扰动结构，检查整条线段留在定义域内，再由微分读出响应与条件数、由精确增量或二阶信息估计余项。

线性运算的余项为零；乘积、平方与二次型的余项是可以精确写出的高次项；求逆的余项由逆差恒等式或Neumann级数给出，界为$\lVert\bm A^{-1}\rVert^3\lVert\bm E\rVert^2/(1-q)$，相对条件数是$\kappa_2(\bm A)$；行列式的绝对变化在奇异处也有二次余项控制，相对条件数只在可逆时有意义，行列式数值小并不意味着求逆病态；对数行列式在对称正定情形下可以借最小特征值明确保证整条线段正定，余项界随到奇异边界的余量缩小而放大。

谱分支没有显式公式，先由隐函数定理在简单特征值处得到可微分支。特征值的变化率是扰动在特征方向上的读数$\bm v^\top\bm E\bm v$，特征向量的变化率以特征值之差为分母，谱隙越小，转向越快。一阶预测可靠的条件是扰动相对谱隙小，而不只是扰动小；重根附近应改为比较隔离的特征子空间。

与已有结论相比，第~\ref{chap:condition-and-spectral-perturbation}章的有限扰动界给出对一切方向、一切步长有效的稳定性保证，第~\ref{chap:solution-error-regularization}章评价方程解与算法的误差；本章补充的是指定方向的响应以及线性预测在有限步长下的精度。使用时要保留定义域、扰动结构、谱隙与相对尺度这些条件，不能用基点处的导数代替整段上的检查。
